\unit{Thuyết trình 212 -- Các tiền lược đồ thương}
\sid{212}
\markboth{Thuyết trình 212}{Thuyết trình 212}
\src{99}
\seminar{Năm thứ 13, 1960/61, số 212\\Tháng 2 năm 1961}
\paperhead{CÁC KỸ THUẬT XÂY DỰNG VÀ CÁC ĐỊNH LÝ TỒN TẠI TRONG HÌNH HỌC ĐẠI SỐ\\
III. CÁC TIỀN LƯỢC ĐỒ THƯƠNG}%
  {Alexander GROTHENDIECK}

\subsection*{Mở đầu}
\addcontentsline{toc}{subsection}{Mở đầu}

Các vấn đề được xét trong exposé này khác với các vấn đề của hai exposé trước
ở chỗ ta tìm cách biểu diễn một số hàm tử hiệp biến, chứ không còn là phản
biến, theo các lược đồ biến thiên. Tuy nhiên, phép lấy thương là thiết yếu
trong nhiều bài toán xây dựng của hình học đại số, kể cả các bài toán trong
Exposé I và II ([1], [2]). Chẳng hạn, câu hỏi về \emph{tính hữu hiệu của một
dữ liệu hạ xuống} trên một $T$-tiền lược đồ $X$, đối với một cấu xạ phẳng
trung thành và quasi-compact $T\longrightarrow S$, tương đương với câu hỏi về
sự tồn tại của một thương của $X$ (thỏa các tính chất hợp lý được xét dưới
đây) theo quan hệ tương đương phẳng trên $X$ do dữ liệu hạ xuống xác định; các
câu hỏi nêu trong [1], A, \S~2 c chắc chắn sẽ được giải quyết đồng thời với
các câu hỏi đặt ra trong đoạn 2 của exposé này. Tương tự, \emph{lược đồ
Picard} của một $S$-lược đồ $X$ (xem định nghĩa trong [2], C, \S~3) có thể
được định nghĩa theo nhiều cách như thương của một số lược đồ khác (các lược
đồ của ước hiệu dụng, hoặc của các phép nhúng vào một lược đồ xạ ảnh) theo
các quan hệ tương đương phẳng; trong khi đó, việc định nghĩa và xây dựng các
lược đồ phụ trợ này đơn giản hơn về mặt kỹ thuật: thật vậy, chúng là các lược
đồ dạng $\underline{\operatorname{Hom}}_S(X,Y)$ và các biến thể được định
nghĩa trong [2], C, \S~2, và việc xây dựng chúng sẽ là nội dung của exposé
tiếp theo (dưới các giả thiết thích hợp về tính xạ ảnh). Vì thế, bằng cách kết
hợp các kết quả của exposé này và exposé tiếp theo, ta thu được phép xây dựng
các lược đồ Picard dưới những giả thiết thích hợp.

Bài toán lấy thương trong phạm trù các tiền lược đồ vẫn còn đặt ra một số câu
hỏi chưa được giải quyết. Câu hỏi quan trọng nhất được nêu trong đoạn 8. Hiện
nay nó là trở ngại duy nhất đối với việc xây dựng \emph{các lược đồ moduli
trên các số nguyên cho các đường cong có giống tùy ý, các đa tạp abel phân
cực}, v.v. Điều đó cho thấy lời giải của nó xứng đáng với nỗ lực của các
chuyên gia về nhóm đại số.
\src{100}
\subsection*{1. Quan hệ tương đương, quan hệ tương đương hữu hiệu}
\addcontentsline{toc}{subsection}{1. Quan hệ tương đương, quan hệ tương đương hữu hiệu}

Cho $\mathcal C$ là một phạm trù và $X$ là một vật của $\mathcal C$. Một cặp
cấu xạ
\[
  p_1,p_2:R\rightrightarrows X
\]
được gọi là một \emph{cặp tương đương} trong $\mathcal C$, có \emph{đích}
$X$ và \emph{nguồn} $R$, nếu với mọi vật $T$ của $\mathcal C$, hai ánh xạ
tương ứng
\[
  p_1(T),p_2(T):R(T)\rightrightarrows X(T)
\]
(trong đó, với mỗi vật $Y$ của $\mathcal C$, ta đặt
$Y(T)=\operatorname{Hom}(T,Y)$) xác định một ánh xạ
\[
  R(T)\longrightarrow X(T)\times X(T)
\]
cảm sinh một song ánh từ $R(T)$ lên đồ thị của một quan hệ tương đương trên
tập $X(T)$. Ta đưa vào một quan hệ tương đương hiển nhiên giữa các cặp tương
đương có đích $X$; một lớp tương đương đối với quan hệ sau này được gọi là
một \emph{quan hệ tương đương của $\mathcal C$ trên $X$}, hoặc đơn giản là
một quan hệ tương đương nếu không thể gây nhầm lẫn.

Nếu $X\times X$ tồn tại, việc cho một quan hệ tương đương trên $X$ tương
đương với việc cho một vật con $R$ của $X\times X$ sao cho, với mọi vật $T$
của $\mathcal C$, tập con của
\[
  (X\times X)(T)=X(T)\times X(T)
\]
tương ứng với $R(T)$ là đồ thị của một quan hệ tương đương trên $X(T)$. Ký
hiệu $p_1$ và $p_2$ là các cấu xạ từ $R$ đến $X$ cảm sinh bởi các phép chiếu
$\operatorname{pr}_1$ và $\operatorname{pr}_2$, điều kiện trên nói rằng
$(p_1,p_2)$ là một cặp tương đương. Ta cũng có thể biểu diễn bằng biểu đồ
trong $\mathcal C$ (với điều kiện tồn tại $X\times X$ và tích thớ
$(R,p_2)\times_X(R,p_1)$) các tiên đề theo nghĩa tập hợp của một quan hệ
tương đương đối với các $R(T)$ trong các $X(T)$, phù hợp với nguyên lý tổng
quát được giải thích trong [2], A, \S~1. Ta sẽ không cần đến điều này.

Mỗi khi có một cặp cấu xạ $(p_1,p_2)$ cùng nguồn $R$ và cùng đích $X$, ta có
thể định nghĩa \emph{đối hạt nhân} của cặp ấy là vật $Y$ của $\mathcal C$
biểu diễn hàm tử hiệp biến theo $Z$
\[
  \operatorname{Hom}_{p_1,p_2}(X,Z),
\]
tức tập các cấu xạ $u$ từ $X$ đến $Z$ sao cho $up_1=up_2$. Nếu $Y$ tồn tại,
nó được xác định duy nhất đến một đẳng cấu duy nhất. Ta ký hiệu nó là
$X/(p_1,p_2)$ hoặc, theo lạm dụng ký hiệu, là $X/R$; ký hiệu sau chủ yếu được
dùng khi
\src{101}$(p_1,p_2)$ là một cặp tương đương: khi đó, theo thông lệ, ta đồng
nhất trong ký hiệu quan hệ tương đương do cặp ấy xác định với $R$. Lưu ý rằng
nếu xem $Y$ là một thương của $X$, thì quả thật nó chỉ phụ thuộc vào quan hệ
tương đương do cặp $(p_1,p_2)$ xác định.

Bây giờ, xét một cấu xạ
\[
  f:X\longrightarrow Y
\]
cho phép xem $X$ như một «vật trên $Y$», và giả sử tích thớ
\[
  R(f)=X\times_Y X
\]
tồn tại. Gọi $p_1$ và $p_2$ là hai phép chiếu của nó. Khi đó $(p_1,p_2)$ là
một cặp tương đương, gọi là cặp \emph{liên kết} với cấu xạ $f$. Do đó, cặp
này xác định một quan hệ tương đương, gọi là quan hệ liên kết với $f$.

Ta nói một cặp cấu xạ $(p_1,p_2)$ có đích $X$ và cùng nguồn $R$ là một
\emph{cặp tương đương hữu hiệu} nếu
\begin{enumerate}
  \item[(i)] đối hạt nhân $X/(p_1,p_2)=Y$ tồn tại;
  \item[(ii)] tích thớ $X\times_Y X$ tồn tại;
  \item[(iii)] cấu xạ $R\longrightarrow X\times_Y X$ có các thành phần $p_1$ và
  $p_2$ là một đẳng cấu.
\end{enumerate}
Khi đó, $(p_1,p_2)$ quả là một cặp tương đương. Ta cũng nói quan hệ tương
đương do nó xác định là một \emph{quan hệ tương đương hữu hiệu}.

Ta nói một cấu xạ $f:X\longrightarrow Y$ là một \emph{toàn cấu hữu hiệu}
nếu
\begin{enumerate}
  \item[(i)] tích thớ $X\times_Y X=R$ tồn tại;
  \item[(ii)] thương $X/(p_1,p_2)$ tồn tại, trong đó $p_1$ và $p_2$ là hai phép
  chiếu từ $R$ đến $X$;
  \item[(iii)] cấu xạ $X/(p_1,p_2)\longrightarrow Y$ do $f$ cảm sinh là một
  đẳng cấu.
\end{enumerate}
Khi đó, $f$ quả là một toàn cấu, và thậm chí là một toàn cấu chặt (xem [1],
A, \S~2.3); điều đảo lại đúng nếu tích thớ $X\times_Y X$ tồn tại. Ta cũng
gọi vật thương của $X$ do toàn cấu $f$ xác định là một \emph{thương hữu hiệu}
của $X$.

Các định nghĩa trên dẫn đến «tương ứng Galois» sau đây:

\src{102}
\result{MỆNH ĐỀ}{1.1} Có một tương ứng song ánh, bảo toàn các thứ tự tự
nhiên, giữa tập các quan hệ tương đương hữu hiệu $R$ trên $X$ và tập các
thương hữu hiệu $Y$ của $X$: quan hệ $R$ ứng với thương hữu hiệu $X/R$, còn
thương $Y$ ứng với quan hệ tương đương hữu hiệu do phép chiếu chính tắc
$X\longrightarrow Y$ xác định (quan hệ này được cho bởi tích thớ $X\times_Y X$
cùng với hai phép chiếu của nó).

Trong các phạm trù rất tốt (các tập hợp, các bó tập hợp, v.v.), mọi thương đều
hữu hiệu và mọi quan hệ tương đương đều hữu hiệu. Điều này không còn đúng
trong những phạm trù như phạm trù các tiền lược đồ trên một tiền lược đồ $S$
cho trước, ngay cả khi $S$ là phổ của một trường, và ngay cả khi chỉ xét các
lược đồ hữu hạn trên $S$. Những vấn đề về tính hữu hiệu, và thậm chí (đối với
các tiền lược đồ không hữu hạn trên $S$) những vấn đề về sự tồn tại của các
thương, thường tỏ ra tế nhị.

\subsection*{2. Ví dụ: các tiền lược đồ hữu hạn trên $S$}
\addcontentsline{toc}{subsection}{2. Ví dụ: các tiền lược đồ hữu hạn trên S}

Cho $\mathcal C$ là phạm trù các tiền lược đồ hữu hạn trên $S$, với giả thiết
cơ sở ấy là Noether địa phương. Phạm trù này tương đương với phạm trù đối của phạm
trù các bó đại số giao hoán coherent trên $S$, hoặc, nếu $S$ là affine có
vành $A$, với phạm trù đối của phạm trù các $A$-đại số hữu hạn trên $A$ (tức
là các môđun hữu hạn sinh trên $A$). Từ đó suy ra ngay rằng trong $\mathcal C$
tồn tại các giới hạn xạ ảnh hữu hạn và các giới hạn quy nạp hữu hạn. Đối với
loại thứ nhất, điều này đã được biết rõ (không cần bất kỳ giả thiết hữu hạn
nào\ldots). Như vậy, tích thớ của các tiền lược đồ $X,Y$ trên $S$ tương ứng
với tích tensor $B\otimes_A C$ của các đại số tương ứng; hạt nhân của hai cấu
xạ $X\rightrightarrows Y$, được xác định bởi hai đồng cấu $A$-đại số
$u,v:C\rightrightarrows B$, tương ứng với thương của $B$ theo iđêan sinh bởi
các phần tử $u(c)-v(c)$, v.v. Đối với các giới hạn quy nạp hữu hạn, chỉ cần
xét, một mặt, các tổng hữu hạn, tương ứng với tích thông thường của các
$A$-đại số, và mặt khác, các đối hạt nhân của các cặp cấu xạ
$X\rightrightarrows Y$; quả thật, chúng tương ứng (như thấy ngay) với vành con
của $C$ gồm những phần tử mà trên đó các đồng cấu
$u,v:C\rightrightarrows B$ trùng nhau (vành con này hữu hạn trên $A$ nhờ giả
thiết Noether). Hơn nữa, lưu ý rằng nhờ giả thiết Noether, ta có thể chứng
minh rằng các giới hạn quy nạp hữu hạn, và đặc biệt là các thương, được xây
dựng như trên trong phạm trù $\mathcal C$ các tiền lược đồ hữu hạn trên $S$,
thực ra là các thương trong phạm trù của \emph{tất cả} các tiền lược đồ.
\src{103}
Như đã nói trong [1], trong $\mathcal C$ có những toàn cấu không hữu hiệu
(hoặc không chặt\footnote{Văn bản in bỏ từ «không» trước «chặt». Trong phạm
trù đang xét, các tích thớ tồn tại, và đoạn trước xác lập chính xác sự tương
đương giữa tính hữu hiệu và tính toàn cấu chặt.}; hai điều này là như nhau vì
các tích thớ tồn tại). \emph{Tôi không biết liệu các quan hệ tương đương có
luôn hữu hiệu hay không} khi không đặt giả thiết nào về tính phẳng. Theo hướng
này, tôi chỉ thu được những kết quả khẳng định rất bộ phận, nhưng không thể
thiếu cho việc chứng minh định lý cơ bản trong lý thuyết hình thức về moduli
(xem [2], B, đl. 1). Lưu ý rằng trong bài toán đã nêu, ta dễ dàng quy về
trường hợp $S$ là phổ của một vành Artin địa phương $A$ có trường thặng dư
đóng đại số. Nhưng ngay cả khi $A$ là một trường, câu trả lời vẫn chưa được
biết.

Ta cũng có thể xét trường hợp một tiền lược đồ $X$ trên $S$ không còn được
giả sử là hữu hạn trên $S$, nhưng được trang bị một quan hệ tương đương $R$
trên $X$ sao cho $p_1:R\longrightarrow X$ là một cấu xạ hữu hạn. Khi đó, ta
nói $R$ là một quan hệ tương đương hữu hạn. Để đơn giản, giả sử $S$ và $X$
đều affine (do đó $R$ affine, nên tình huống được quy về một tình huống đại số
giao hoán thuần túy); \emph{ngay cả trong trường hợp này, ta cũng chưa biết có
tồn tại thương $X/R=Y$ hay không và cấu xạ chính tắc
$X\longrightarrow Y$ có hữu hạn hay không}. (Trường hợp đơn giản nhất là khi
$S$ là phổ của một trường $k$, còn $X$ là phổ của $k[t]$, tức đường thẳng
affine.) Dĩ nhiên, nếu hai bài toán trên đều có lời giải khẳng định, thì trong
tình huống hiện tại ta có thể kết luận rằng $R$ là hữu hiệu. Lưu ý rằng bài
toán về \emph{sự tồn tại} của một thương $Y$ và tính hữu hạn của
$f:X\longrightarrow Y$ được đặt ra hoàn toàn theo cùng một cách nếu, thay cho
đồ thị của một quan hệ tương đương trên $X$, ta chỉ có một tiền đồ thị tương
đương trên $X$, theo nghĩa của đoạn 4.

Vấn đề lấy thương theo một quan hệ tương đương hữu hạn ít nhiều tùy ý xuất
hiện khi dựng các tiền lược đồ bằng cách «dán» những tiền lược đồ $X_i$ đã cho
dọc theo một số tiền lược đồ con đóng; luật dán được biểu thị chính xác bởi
một quan hệ tương đương hữu hạn trên tiền lược đồ tổng $X$ của các $X_i$. Ta
cũng nên chờ đợi rằng lời giải cho các bài toán đặt ra ở đây và cho những biến
thể khác nhau sẽ là một điều kiện sơ bộ để xây dựng một kỹ thuật tổng quát về
các phép dựng không xạ ảnh, theo hướng được mở ra trong [2].

Kết quả tổng quát theo chiều khẳng định duy nhất mà người biên soạn biết là
kết quả sau:

\result{MỆNH ĐỀ}{2.1} Cho $S$ là một tiền lược đồ Noether địa phương, $s$ là
một điểm của $S$, và $\Omega$ là một mở rộng trường đóng đại số của $k(s)$.
Xét
\src{104}«hàm tử thớ» tương ứng $F$, gán cho mọi $S$-lược đồ $X$ hữu hạn
trên $S$ tập các điểm của $X_s$ nhận giá trị trong $\Omega$. Hàm tử này (hiển
nhiên khớp trái) là \emph{khớp phải}, tức giao hoán với các giới hạn quy nạp
hữu hạn, và đặc biệt với các đối hạt nhân của những cặp cấu xạ.

Áp dụng kết quả này cho mọi «điểm hình học» của $S$, ta suy ra rằng phạm trù
«thương» $\mathcal C'$ của $\mathcal C$, thu được bằng cách lập luận «modulo
các cấu xạ radiciel toàn ánh» (tức là ghép thêm một cách hình thức các nghịch
đảo của những cấu xạ đang xét), là một phạm trù «hình học», nghĩa là thỏa mãn
cùng những tính chất «có bản chất hữu hạn» như phạm trù các tập hợp. Đặc biệt,
mọi quan hệ tương đương trong đó đều hữu hiệu. Điều này kéo theo rằng nếu $R$
là một quan hệ tương đương trên một $X$ hữu hạn trên $S$, thì cấu xạ chính tắc
\[
  R\longrightarrow X\times_Y X\qquad (\text{trong đó }Y=X/R)
\]
là \emph{radiciel và toàn ánh} (thực ra là một phép nhúng đóng toàn ánh, vì nó
là một đơn cấu).

\subsection*{3. Trường hợp một nhóm tác tử}
\addcontentsline{toc}{subsection}{3. Trường hợp một nhóm tác tử}

Lại giả sử $\mathcal C$ là một phạm trù tùy ý. Cho $G,X$ là các vật của
$\mathcal C$, và giả sử $G$ là một $\mathcal C$-nhóm tác tử trên vật $X$ của
$\mathcal C$. Điều này có nghĩa (xem [2], A, \S~1) rằng với mọi vật $T$ của
$\mathcal C$, ta đã cho một cấu trúc nhóm trên $G(T)$ và một cấu trúc tập hợp
có nhóm tác tử $G(T)$ trên $X(T)$, sao cho khi $T$ biến thiên, những cấu trúc
đang xét «biến thiên theo hàm tử» theo $T$. Nếu trong $\mathcal C$ tồn tại các
tích $G\times G$ và $G\times X$, thì một cấu trúc như vậy cũng có thể được
xác định bởi một cặp cấu xạ
\[
  G\times G\longrightarrow G,
  \qquad \pi:G\times X\longrightarrow X,
\]
thỏa điều kiện rằng với mọi vật $T$ của $\mathcal C$, các luật hợp thành tương
ứng trên các tập $G(T)$ và $X(T)$ làm cho $G(T)$ thành một nhóm tác động trên
$X(T)$. Việc diễn đạt tiên đề này bằng tính giao hoán của một số biểu đồ trong
$\mathcal C$ là dễ nhưng dài dòng, và thực ra hoàn toàn vô ích trong mọi
trường hợp mà tôi biết.

Giả sử $G\times X$ tồn tại, và xét hai cấu xạ
\[
  p_1,p_2:G\times X\rightrightarrows X
\]
được cho bởi\footnote{Văn bản in ghi \texttt{pr\_1}. Vì cả hai cấu xạ đều có
đối đích $X$, cấu xạ thứ nhất phải là \texttt{pr\_2}; ở đây ta sửa lỗi in.}

\src{105}
\[
  p_1=\operatorname{pr}_2,\qquad p_2=\pi.
\]
Ta thấy ngay rằng $(p_1,p_2)$ là một cặp tương đương khi và chỉ khi, với mọi
vật $T$ của $\mathcal C$, ánh xạ
\[
  G(T)\times X(T)\simeq(G\times X)(T)
  \longrightarrow X(T)\times X(T),
\]
do cặp này xác định là đơn ánh, tức là khi nhóm $G(T)$ tác động
\emph{tự do} trên tập $X(T)$, nghĩa là khi $g\in G(T)$, $x\in X(T)$ và
$g\cdot x=x$ kéo theo $g$ là phần tử đơn vị của nhóm $G(T)$. Khi đó, ta nói
rằng $G$ \emph{tác động tự do trên $X$} (hoặc $X$ là một
\emph{$\mathcal C$-không gian chính dưới $G$}). Quan hệ tương đương liên kết
với cặp $(p_1,p_2)$ khi ấy được gọi là \emph{quan hệ tương đương do nhóm $G$
tác động tự do trên $X$ xác định}. Khi $X\times X$ cũng tồn tại và ta xét cấu
xạ
\[
  p:G\times X\longrightarrow X\times X
\]
do cặp $(p_1,p_2)$ xác định, điều kiện $G$ tác động tự do có nghĩa là $p$ là
một \emph{đơn cấu}.

Dĩ nhiên, ngay cả khi $G$ không tác động tự do trên $X$, ta vẫn muốn có các
tiêu chuẩn về sự tồn tại của một thương của $X$ theo $G$, tức đối hạt nhân của
cặp $(p_1,p_2)$ trên.

Đối hạt nhân đang xét thường được ký hiệu là $X/G$, hoặc tốt hơn là
$G\backslash X$ nếu $G$ tác động bên trái (ký hiệu trước được dành cho trường
hợp $G$ tác động bên phải). Lưu ý rằng ngay cả khi «ảnh» của $G\times X$ qua
$p$ tồn tại (chẳng hạn ảnh này được định nghĩa là vật con nhỏ nhất của
$X\times X$ mà qua đó $p$ phân tích được), ký hiệu là $R$, thì nói chung nó
không phải là một quan hệ tương đương trên $X$. Nếu khi ấy ta cố lấy trực tiếp
thương theo $R$ (hay chính xác hơn, theo cặp cấu xạ từ $R$ đến $X$ cảm sinh bởi
hai phép chiếu $\operatorname{pr}_i$), thì ta làm mất không thể khôi phục
những đặc trưng riêng của cặp ban đầu $(p_1,p_2)$. Vì vậy, điều quan trọng là
tìm một sự khái quát hóa của khái niệm quan hệ tương đương, áp dụng trực tiếp
cho cặp được xác định bởi một $\mathcal C$-nhóm tác tử.

\subsection*{4. Tiền quan hệ tương đương}
\addcontentsline{toc}{subsection}{4. Tiền quan hệ tương đương}

Nhắc lại rằng một phạm trù trong đó mọi cấu xạ đều là đẳng cấu được gọi là một
\emph{groupoid}. Mặt khác, một phạm trù phải được định nghĩa là gồm hai tập nền
$(X,R)$, tập các \emph{vật} và tập các \emph{mũi tên}, được trang bị những cấu
trúc sau đây:
\src{106}
\begin{enumerate}
\item[(i)] Một cặp ánh xạ
\[
  p_1,p_2:R\rightrightarrows X
\]
được gọi là \emph{ánh xạ nguồn} và \emph{ánh xạ đích},

\item[(ii)] Một ánh xạ
\[
  \pi:(R,p_2)\times_X(R,p_1)\longrightarrow R
\]
được gọi là \emph{ánh xạ hợp thành}.
\end{enumerate}
Các dữ liệu này phải thỏa mãn những tiên đề quen thuộc, mà ta sẽ không nhắc
lại ở đây, và có thể được diễn đạt bằng tính giao hoán của một số biểu đồ,
cùng với sự tồn tại của một ánh xạ $D:X\to R$ (nhất thiết duy nhất) làm cho
hai biểu đồ khác giao hoán. Ánh xạ $D$ gửi một vật tới mũi tên đồng nhất tương ứng
và thỏa mãn
\[
  p_1D=p_2D=\operatorname{id}_X.
\]
Việc nói phạm trù là một groupoid khi đó có nghĩa là tồn tại một ánh xạ
\[
  s:R\longrightarrow R
\]
(nhất thiết duy nhất), gọi là \emph{phép lấy nghịch đảo} của $R$, biến mỗi
mũi tên thành một mũi tên nghịch đảo. Điều này có thể được diễn đạt bằng tính
giao hoán của bốn biểu đồ khác được tạo từ $s$, $\Delta$ và các dữ liệu trước
đó; hai biểu đồ đầu là
\[
  p_1s=p_2,\qquad p_2s=p_1.
\]

Nhắc lại các điểm này, các định nghĩa tổng quát trong ([2], A, \S~1) đặc biệt
cho thấy thế nào là một cấu trúc $\mathcal C$-\emph{phạm trù}, tương ứng một
cấu trúc $\mathcal C$-\emph{groupoid}, trên một cặp vật $X,R$ của một phạm trù
tùy ý $\mathcal C$: theo định nghĩa, đó là dữ liệu, với mỗi vật $T$ của
$\mathcal C$, của một cấu trúc phạm trù, tương ứng một cấu trúc groupoid, theo
nghĩa tập hợp, có tập các vật là $X(T)$ và tập các mũi tên là $R(T)$, các cấu
trúc này \guillemotleft biến thiên theo hàm tử\guillemotright{} khi $T$ biến
thiên. Do đó, điều này xác định hai cấu xạ
\[
  p_1,p_2:R\rightrightarrows X
\]
được gọi là \emph{cấu xạ nguồn} và \emph{cấu xạ đích}, và khi tích thớ cần
thiết tồn tại, một cấu xạ
\src{107}
\[
  \pi:(R,p_2)\times_X(R,p_1)\longrightarrow R
\]
được gọi là \emph{cấu xạ hợp thành}. Ba dữ liệu này đủ để xác định cấu trúc
phạm trù trên $X,R$; tiên đề đặt lên các dữ liệu ấy là như sau: với mọi $T$,
ba dữ liệu tương ứng cho $X(T),R(T)$ xác định trên cặp tập hợp này một cấu
trúc phạm trù (tương ứng, một cấu trúc groupoid). Khi thích hợp, điều này có
thể được diễn đạt bằng tính giao hoán của một số biểu đồ có chứa một cấu xạ
được xác định duy nhất
\[
  D:X\longrightarrow R
\]
và, trong trường hợp các groupoid, một cấu xạ được xác định duy nhất
\[
  s:R\longrightarrow R;
\]
các biểu đồ được viết tường minh đúng như trong trường hợp «theo nghĩa tập
hợp». May thay, cách diễn giải dài dòng này của các tiên đề là vô ích trong
thực hành. Lợi ích lý thuyết duy nhất của khả năng biểu diễn dữ liệu và tiên
đề bằng các cấu xạ giữa một số tích thớ và các đẳng thức giữa những cấu xạ ấy
là như sau: nếu $F:\mathcal C\to\mathcal C'$ là một hàm tử khớp trái (tức giao
hoán với các tích hữu hạn và các tích thớ), thì nó biến một
$\mathcal C$-phạm trù (tương ứng, một $\mathcal C$-groupoid) thành một
$\mathcal C'$-phạm trù (tương ứng, một $\mathcal C'$-groupoid), với điều kiện
các tích hữu hạn và các tích thớ tồn tại trong $\mathcal C$.

Trong thực hành, điều quan trọng là biết diễn giải các cấu xạ
$p_1,p_2,\pi,D,s$ như những \emph{phép toán đơn hình} trong một vật bán đơn
hình hoặc đơn hình thích hợp của $\mathcal C$, ít nhất khi các tích thớ tồn
tại trong $\mathcal C$. Để cố định thuật ngữ, ta đưa vào \emph{phạm trù
$\mathcal S$ các đơn hình chuẩn} là phạm trù mà các vật là những tập hữu hạn
có dạng
\[
  \Delta_n=[0,n],\qquad\text{trong đó }n\geq 0
\]
(khoảng các số nguyên từ $0$ đến $n$), và các cấu xạ hay mũi tên là
\emph{các ánh xạ tùy ý} giữa những tập hữu hạn ấy. Lưu ý rằng phạm trù
$\mathcal S$ tương đương với phạm trù các tập hữu hạn \emph{khác rỗng}, trong
đó các cấu xạ là những ánh xạ giữa các tập hữu hạn. Trong $\mathcal S$, đối
tích của một họ hữu hạn \emph{khác rỗng} các vật hiển nhiên tồn tại, cũng như
phép đẩy ra của hai vật dưới một vật thứ ba (phép toán đối ngẫu với tích thớ).
Ký hiệu $\mathcal S'$ là phạm trù con của $\mathcal S$ có cùng các vật, nhưng
trong đó các cấu xạ là \emph{các ánh xạ tăng} giữa các $\Delta_n$. Phạm trù
này tương đương với phạm trù các tập hữu hạn khác rỗng được sắp thứ tự toàn
phần. Trong phạm trù này,
\src{108}
đối tích của hai vật không bao giờ tồn tại, và phép đẩy ra của hai vật $A,B$
dưới một vật thứ ba $C$ nói chung không tồn tại (chẳng hạn, lấy
$C=\Delta_0$, $A=B=\Delta_1$, với hai ánh xạ cấu trúc $u:C\to A$ và
$v:C\to B$ bằng nhau). Tuy nhiên, nó tồn tại trong một số trường hợp; chẳng
hạn
\[
  C=\Delta_0,\qquad A=\Delta_m,\qquad B=\Delta_n,\qquad
  u(0)=m,\qquad v(0)=0,
\]
và trong trường hợp này
\[
  \Delta_m\amalg_{\Delta_0}\Delta_n=\Delta_{m+n}.
\]

Một \emph{vật đơn hình} (tương ứng, một \emph{vật bán đơn hình}) trong một
phạm trù $\mathcal C$ theo định nghĩa là một hàm tử phản biến $K$ từ
$\mathcal S$ (tương ứng từ $\mathcal S'$) vào $\mathcal C$. Do đó, một vật
đơn hình xác định một vật bán đơn hình bằng phép hạn chế, nhưng khái niệm thứ
nhất khác khái niệm thứ hai về căn bản bởi sự hiện diện của \emph{các phép
toán đối xứng} trên các $K_n=K(\Delta_n)$, tương ứng với ảnh qua hàm tử $K$
của các phần tử thuộc nhóm đối xứng trên $n+1$ phần tử (được xem là nhóm các
tự đẳng cấu của $\Delta_n$ trong $\mathcal S$).

Với mọi $n$, gọi $\Delta'_n$ (tương ứng $\Delta''_n$) là phạm trù hữu hạn có
tập các vật là $\Delta_n$, được xác định bởi quan hệ thứ tự hỗn loạn (tương
ứng, thứ tự toàn phần tự nhiên) trên $\Delta_n$ (tập các mũi tên là đồ thị của
quan hệ thứ tự ấy). Hiển nhiên $\Delta'_n$ (tương ứng $\Delta''_n$) phụ thuộc
theo hàm tử vào vật $\Delta_n$ của $\mathcal S$ (tương ứng $\mathcal S'$).
Nếu $Z$ là một phạm trù, thì $\operatorname{Hom}(\Delta'_n,Z)$ (tương ứng
$\operatorname{Hom}(\Delta''_n,Z)$), khi $\Delta_n$ biến thiên, là một hàm tử
từ phạm trù $\mathcal S$ (tương ứng $\mathcal S'$) vào phạm trù các tập hợp,
tức một tập đơn hình (tương ứng, một tập bán đơn hình), được gọi là
\emph{liên kết với phạm trù $Z$}; ký hiệu chúng lần lượt là $Z'$ và $Z''$.
Hơn nữa, có một đồng cấu tự nhiên hiển nhiên từ tập bán đơn hình liên kết
với $Z'$ vào
$Z''$, và đồng cấu này là một đẳng cấu khi và chỉ khi $Z$ là một groupoid. Với
các quy ước này:

\result{MỆNH ĐỀ}{4.1} Hàm tử $Z\longmapsto Z''$ từ phạm trù các phạm trù vào
phạm trù các tập bán đơn hình là đầy và trung thành, và xác định một tương
đương giữa phạm trù các \emph{phạm trù} với phạm trù các tập bán đơn hình,
tức các hàm tử phản biến $K$ từ $\mathcal S'$ vào $(\mathrm{Set})$, biến các
phép đẩy ra $A\amalg_C B$ thuộc kiểu được chỉ ra ở trên thành các tích thớ của
các tập hợp. Tương tự, hàm tử $Z\longmapsto Z'$ từ phạm trù các groupoid vào
phạm trù các tập đơn hình là đầy và trung thành, và xác định một tương
đương giữa phạm trù các \emph{groupoid} với phạm trù các
\src{109}
các tập đơn hình, tức các hàm tử phản biến $K$ từ $\mathcal S$ vào
$(\mathrm{Ens})$, biến các phép đẩy ra thành các tích thớ.

Do đó, có thể xem các phạm trù như những tập bán đơn hình đặc biệt, và các
groupoid như những tập đơn hình đặc biệt, dĩ nhiên với điều kiện phải lập luận
«sai khác một đẳng cấu», như bắt buộc khi diễn giải một số cấu trúc bằng những
cấu trúc khác. Phương pháp thông thường quy về trường hợp lý thuyết tập hợp
khi đó kéo theo:

\result{HỆ QUẢ}{4.2} Mệnh đề trên vẫn đúng khi thay các phạm trù, các
groupoid và các tập đơn hình bằng các $\mathcal C$-phạm trù,
$\mathcal C$-groupoid và các $\mathcal C$-vật đơn hình, với điều kiện các tích
thớ tồn tại trong $\mathcal C$.

Vật bán đơn hình $K$ trong $\mathcal C$ liên kết với một phạm trù
$(X,R,\ldots)$ trong $\mathcal C$ có thể được viết tường minh bằng cách diễn
giải thành phần $K_n=K(\Delta_n)$ của $K$ là tích thớ lặp $n$ lần của
$(R,p_1)$ trên $X$, hoặc tốt hơn, bằng công thức truy hồi
\[
  K_0=X,\qquad
  K_n=(K_{n-1},p_{n-1}^{(n-1)})\times_X(R,p_1),
\]
trong đó các $p_i^{(n-1)}$ ($0\leq i\leq n-1$) là các phép chiếu tự nhiên từ
$K_{n-1}$ xuống $X$ (cũng được định nghĩa bằng truy hồi). Theo cách này,
$p_1,p_2,\pi,D,s$ được diễn giải là các phép toán đơn hình lần lượt tương ứng
với các cấu xạ sau trong $\mathcal S$: mặt $0$ của $\Delta_1$, mặt $1$ của
$\Delta_1$, mặt $(0,2)$ của $\Delta_2$, phép suy biến
$\Delta_1\to\Delta_0$, và phép đối xứng của $\Delta_1$. Mọi phép toán bán đơn
hình khác (tương ứng, đơn hình) đều được suy ra một cách hình thức từ bốn
phép toán trước (tương ứng, năm phép toán) bằng phép hợp thành và các tích thớ.

Nay gọi \emph{tiền quan hệ tương đương} trên một vật $X$ của một phạm trù là
việc cho một groupoid mà vật của các vật (nếu có thể nói vậy) là $X$. Vì thế,
dữ liệu ấy bao gồm, ngoài những thứ khác, một vật $R$ và hai cấu xạ
\[
  p_1,p_2:R\rightrightarrows X.
\]
Tuy nhiên, lưu ý rằng riêng các dữ liệu này không xác định cấu trúc đang xét,
khác với trường hợp các cặp tương đương. Trong exposé này, ta quan tâm đến
khái niệm ấy nhằm thu được các tiêu chuẩn để có thể lấy thương, tức là lập đối
hạt nhân của cặp $(p_1,p_2)$. Vì vậy, phát biểu của bài toán không dùng đến
các dữ liệu bổ sung vốn có trong một groupoid. Trong việc chứng minh các kết
quả sẽ nêu, ta sẽ
\src{110}
vẫn sử dụng các dữ liệu bổ sung ấy, và đặc biệt là các phép toán đơn hình
(kể cả các phép toán đối xứng) cho đến chiều $3$ (có liên quan đến tích thớ
lặp bốn lần của $R$ trên $X$).

Một quan hệ tương đương trên một vật $X$ của $\mathcal C$ xác định một tiền
quan hệ tương đương: chỉ cần xét điều này trong trường hợp tập hợp, và khi ấy,
với một quan hệ tương đương trên một tập $X$, ta liên kết groupoid có tập vật
là $X$, còn tập mũi tên là đồ thị của quan hệ tương đương.

Một $\mathcal C$-monoid $G$ tác động lên một vật $X$ của $\mathcal C$ xác định
một $\mathcal C$-phạm trù có các vật cơ sở là $R=G\times X$ và $X$ (với điều
kiện $G\times X$ tồn tại), và là một $\mathcal C$-groupoid khi và chỉ khi $G$
là một nhóm. Một lần nữa, chỉ cần kiểm tra điều này trong trường hợp tập hợp.
Khi đó, phép hợp thành các mũi tên $(g,a)$ và $(g',g\cdot a)$ được định nghĩa
bởi
\[
  (g',g\cdot a)(g,a)=(g'g,a),
\]
tức là, nếu $a,b\in X$, thì theo định nghĩa $\operatorname{Hom}(a,b)$ là tập
chuyển từ $a$ đến $b$, và các cấu xạ được hợp thành nhờ phép hợp thành các phần
tử của $G$.

\medskip
\noindent\textbf{NHẬN XÉT. ---} Có thể tránh các khó khăn logic do một phát
biểu như (4.1) gây ra bằng cách hiểu ngầm rằng mọi vật đang xét đều nằm trong
một «Vũ trụ» cố định (bản thân nó là một tập hợp).

\fsection{5}{Thương theo một quan hệ tương đương hữu hạn và phẳng}

\result{ĐỊNH LÝ}{5.1} Cho $X=\operatorname{Spec}(B)$ là một lược đồ affine và
$R$ là một tiền quan hệ tương đương trên $X$, mà thành phần $R_1$ là affine,
đặt $R_1=\operatorname{Spec}(C)$. Giả sử phép chiếu thứ nhất
$p_1:R_1\to X$ là một cấu xạ hữu hạn và tự do địa phương; nói cách khác, giả
sử đồng cấu vành tương ứng $p_1^*:B\to C$ làm cho $C$ trở thành một
$B$-môđun xạ ảnh hữu hạn sinh. Gọi $A$ là vành con của $B$, là hạt nhân của
cặp đồng cấu
\[
  p_1^*,p_2^*:B\rightrightarrows C
\]
(= tập hợp các phần tử $b$ sao cho $p_1^*(b)=p_2^*(b)$). Đặt
$Y=\operatorname{Spec}(A)$, và gọi $f:X\to Y$ là cấu xạ được xác định bởi phép
nhúng $A$ vào $B$. Trong các điều kiện này:
\begin{enumerate}
\item[(i)] $B$ nguyên trên $A$, tức là $f$ là một cấu xạ nguyên.

\item[(ii)] Cấu xạ $f$ là toàn ánh, và các thớ của nó là các lớp tương đương
\src{111}
theo nghĩa tập hợp $p_2\allowbreak(p_1^{-1}(x))$ trong $X$ modulo $R$, và
tôpô của $Y$ là tôpô thương của tôpô trên $X$.

\item[(iii)] $Y$ là thương của $X$ theo $R$ trong phạm trù các tiền lược đồ.

\item[(iv)] Nếu $R$ xuất phát từ một \emph{quan hệ tương đương}, thì cấu xạ
$f:X\to Y$ là hữu hạn và tự do địa phương (nghĩa là $B$ là một $A$-môđun xạ
ảnh hữu hạn sinh), và quan hệ tương đương là hữu hiệu, nghĩa là
$R_1\to X\times_Y X$ là một đẳng cấu.
\end{enumerate}

Định lý này tổng quát hóa định lý quen thuộc về trường hợp một nhóm hữu hạn
$G$ tác động bằng các tự đẳng cấu lên vành $B$ và về vành bất biến $A$; chứng
minh của nó tương tự chứng minh quen thuộc. Có thể phát biểu chính xác hơn
(iii) như sau:

\result{HỆ QUẢ}{5.2} Cấu xạ chính tắc
$R_1\to X\times_Y X$ là \emph{toàn ánh}.

Vẫn cho $R$ là một tiền quan hệ tương đương «hữu hạn và tự do địa phương»
trên $X$, nhưng bây giờ $X$ là một tiền lược đồ tùy ý. Giả sử có thể tìm được
một tiền lược đồ $Y$ và một cấu xạ $f:X\to Y$ sao cho $fp_1=fp_2$, đồng thời
dãy các đồng cấu giữa các bó vành trên $Y$
\[
  \mathcal O_Y\longrightarrow f_*(\mathcal O_X)
  \rightrightarrows g_*(\mathcal O_R)
\]
là khớp (trong đó $g=fp_i$). Khi đó, từ định lý suy ra các kết luận tương tự
(i) đến (iv) của định lý; đặc biệt, theo (iii), $Y$ là thương của $X$ theo
$R$, và do đó được xác định sai khác một đẳng cấu duy nhất. Trong các điều
kiện ấy, tiền quan hệ tương đương $R$ trên $X$ được gọi là \emph{khả chấp}.
Với định nghĩa này:

\result{ĐỊNH LÝ}{5.3} Cho $X$ là một tiền lược đồ và $R$ là một tiền quan hệ
tương đương trên $X$ sao cho $p_1:R_1\to X$ là một cấu xạ hữu hạn và tự do địa
phương. Khi đó, $R$ khả chấp khi và chỉ khi mọi lớp tương đương theo nghĩa tập
hợp $p_2(p_1^{-1}(x))$ trong $X$ modulo $R$ đều nằm trong một tập mở affine
(điều kiện này luôn được thỏa mãn nếu mọi tập con hữu hạn của $X$ đều nằm
trong một tập mở affine, chẳng hạn nếu $X$ là tựa xạ ảnh trên một lược đồ
affine).

Thật vậy, không khó để chứng minh rằng khi ấy mọi lớp tương đương modulo $R$
trong $X$ đều nằm trong một tập mở affine \emph{ổn định theo $R$}, và thương
$Y$ được dựng bằng cách dán các mảnh thu được nhờ Định lý 5.1.
\src{112}

\result{HỆ QUẢ}{5.4} Giả sử điều kiện ấy được thỏa mãn và ngoài ra $R$ xuất
phát từ một quan hệ tương đương. Khi đó, quan hệ này là hữu hiệu, nghĩa là
$R_1\to X\times_Y X$ là một đẳng cấu, và $f:X\to Y$ là một cấu xạ hữu hạn và
tự do địa phương.

Từ đó, bằng «hạ xuống», ta suy ra ngay:

\result{HỆ QUẢ}{5.5} Dưới các điều kiện của (5.4), $X$ có hạng $n$ tại mọi
điểm trên $Y$ khi và chỉ khi $(R_1,p_1)$ có hạng $n$ tại mọi điểm trên $X$.
Nếu $X,R_1$ là các $Z$-tiền lược đồ và $p_1,p_2$ là các $Z$-cấu xạ, nên $Y$
là một $Z$-tiền lược đồ, thì $X$ phẳng trên $Z$ khi và chỉ khi $Y$ cũng vậy.

Tóm lại:

\result{CHÚ GIẢI}{} Việc cho một cấu xạ hữu hạn, tự do địa phương và toàn ánh
$f:X\to Y$ giữa các tiền lược đồ tương đương với việc cho một tiền lược đồ
$X$ được trang bị một quan hệ tương đương $R$ sao cho $p_1:R\to X$ là hữu hạn
và tự do địa phương, và mọi lớp $p_2(p_1^{-1}(x))$ đều nằm trong một tập mở
affine.

\medskip
\noindent\textbf{CÁC NHẬN XÉT 5.6.}

\smallskip
\noindent $1^\circ$ Không cần đặt giả thiết Noether nào.

\smallskip
\noindent $2^\circ$ Khái niệm lấy thương này bao gồm, như một trường hợp đặc
biệt, «hạ xuống không khả phân» của CARTIER, tương ứng với việc xác định các
cấu xạ hữu hạn và tự do địa phương $f:X\to Y$ sao cho
$f_*(\mathcal O_X)$ có một $p$-cơ sở đối với $\mathcal O_Y$ ($X$ là một tiền
lược đồ đã cho mà bó $\mathcal O_X$ bị số nguyên tố $p>0$ triệt tiêu). Lưu ý
rằng kết quả được phát biểu dễ dàng mà không cần giả thiết chính quy nào về
các vành địa phương và không cần giả sử $X$ là một lược đồ đại số trên một
trường. Lý thuyết JACOBSON--BOURBAKI thu được bằng cách lấy $X$ là phổ của
một trường có đặc số $p$.

\smallskip
\noindent $3^\circ$ Trước đó, GABRIEL đã thu được một trường hợp riêng của
Định lý (5.3) trong lý thuyết lấy thương đối với các nhóm giao hoán hữu hạn
trên một trường $k$. [So sánh (7.4)].

\fsection{6}{Thương theo một quan hệ tương đương riêng và phẳng}

\result{ĐỊNH LÝ}{6.1} Cho $S$ là một tiền lược đồ Noether địa phương, $X$ là
một $S$-lược đồ tựa xạ ảnh, và $R$ là một tiền quan hệ tương đương trên $X$
sao cho:
\src{113}
\begin{enumerate}
\item[(a)] $p_1:R_1\to X$ là riêng và phẳng;
\item[(b)] $R_1\to X\times_S X$ là một cấu xạ hữu hạn, hoặc tương đương theo
(a), có các thớ hữu hạn (điều kiện tự động được thỏa mãn nếu $R$ xuất phát từ
một quan hệ tương đương).
\end{enumerate}
Dưới các điều kiện này:
\begin{enumerate}
\item[(i)] $X/R=Y$ tồn tại và (nếu $S$ là Noether và nếu $R$ xuất phát từ một
quan hệ tương đương) là tựa xạ ảnh trên $S$.\footnote{Văn bản in không chứa
giả thiết thứ hai. Bản đính chính gắn với
Exposé 212 nói rằng tính tựa xạ ảnh khi ấy dường như chỉ được chứng minh khi
$R$ xuất phát từ một quan hệ tương đương.}

\item[(ii)] Cấu xạ chính tắc $f:X\to Y$ là toàn ánh, riêng và mở; các thớ của
nó là các lớp tương đương $p_2(p_1^{-1}(x))$ trong $X$ modulo $R$. Do đó,
$Y$ đồng nhất với không gian tôpô thương của $X$ theo quan hệ tương đương theo
nghĩa tập hợp do $R$ xác định. Cuối cùng,
$R_1\to X\times_Y X$ là toàn ánh.

\item[(iii)] Nếu $R$ xuất phát từ một quan hệ tương đương, thì quan hệ ấy là
hữu hiệu, nghĩa là $R_1\to X\times_Y X$ là một đẳng cấu, và hơn nữa
$f:X\to Y$ là phẳng (do đó phẳng trung thành).
\end{enumerate}

Để chứng minh, ta quy về (5.1) nhờ các tựa tiết diện thích hợp của $X$ đối với
$R$; chứng minh tương tự việc dựng các nhóm đại số thương trong Séminaire
Chevalley.

Tóm lại:

\result{CHÚ GIẢI}{} Cho $X$ tựa xạ ảnh trên $S$ Noether địa phương. Việc cho
một cấu xạ riêng, phẳng trung thành và toàn ánh $f:X\to Y$ từ $X$ vào một
$S$-tiền lược đồ $Y$ tương đương với việc cho một quan hệ tương đương $R$ trên
$X$ sao cho $p_1:R\to X$ là riêng và phẳng.

Phương pháp ấy cũng cho kết quả sau:

\result{ĐỊNH LÝ}{6.2} Cho $S$ là một tiền lược đồ Noether, $X$ là một tiền
lược đồ thuộc kiểu hữu hạn trên $S$, và $R$ là một tiền quan hệ tương đương trên
$S$-tiền lược đồ $X$; giả sử rằng:
\begin{enumerate}
\item[(a)] $p_1:R_1\to X$ là phẳng và thuộc kiểu hữu hạn;
\item[(b)] cấu xạ $R_1\to X\times_S X$ là tựa hữu hạn (nghĩa là có các thớ
hữu hạn).
\end{enumerate}
Khi đó, tồn tại một tập mở $U$ trù mật trong $X$, \emph{bão hòa} đối với $R$,
sao cho:
\begin{enumerate}
\item[(i)] Nếu $R_U$ là tiền quan hệ tương đương do $R$ cảm sinh trên $U$,
thì $U/R_U$ tồn tại và thuộc kiểu hữu hạn trên $S$.

\item[(ii)] Cấu xạ chính tắc $U\to U/R_U$ là toàn ánh và mở; các thớ của nó
là các lớp tương đương theo nghĩa tập hợp đối với $R_U$ (do đó $U/R_U$ là
không gian tôpô thương của $U$ theo quan hệ tương đương theo nghĩa tập hợp do
$R_U$ xác định). Cuối cùng, cấu xạ
\[
  (R_U)_1\longrightarrow U\times_{U/R_U}U
\]
là toàn ánh.
\src{114}

\item[(iii)] Nếu $R$ xuất phát từ một quan hệ tương đương, thì có thể giả sử
$U\to U/R_U$ phẳng trung thành và $R_U$ hữu hiệu.
\end{enumerate}

Đây là một kết quả có bản chất chủ yếu là «song hữu tỉ».

\medskip
\noindent\textbf{CÁC NHẬN XÉT 6.3.}

\smallskip
\noindent $1^\circ$ Tôi không biết liệu trong (6.1) và (6.2), giả thiết (b) có
thừa hay không. Khi lấy thương theo các nhóm, giả thiết ấy thực tế buộc ta
giới hạn vào trường hợp mọi nhóm ổn định đều là nhóm hữu hạn.

\smallskip
\noindent $2^\circ$ Có thể hỏi liệu có những kết quả tương tự (6.1) và (6.2)
mà không cần giả thiết phẳng hay không. Tôi không biết phản ví dụ nào theo
hướng này. Mặt khác, ngay cả khi vẫn giữ giả thiết phẳng và chỉ xét các quan
hệ tương đương sao cho $p_1:R\to X$ phẳng và tựa hữu hạn (nhưng không hữu
hạn), với $X$ affine, vẫn có thể xảy ra việc $R$ không hữu hiệu: hãy lấy các
quan hệ tương đương cảm sinh trên các tập mở affine phủ đa tạp Nagata (trên đó
một nhóm gồm hai phần tử tác động theo cách «không khả chấp»).

\fsection{7}{Ứng dụng}

Như đã nói trong phần mở đầu, ứng dụng quan trọng nhất của (6.1) là việc dựng
các lược đồ Picard, cũng như việc tìm nghiệm cho nhiều bài toán moduli khác;
chúng ta sẽ trở lại các bài toán này sau.

Ta thu được một chứng minh đơn giản cho kết quả sau đây của SHIMURA:

\result{MỆNH ĐỀ}{7.1} Cho $A$ là một lược đồ abel được xác định trên một vành
định giá rời rạc $V$ có trường phân thức $K$. Khi đó, mọi lược đồ abel $B'$
trên $K$ đẳng giống với một thương của $A\otimes_V K$ đều «có giảm tốt trên
$V$», tức là đẳng cấu với một $B\otimes_V K$,\footnote{Văn bản in ghi
$B\otimes_V A$. Vì $B$ là một $V$-lược đồ và khẳng định liên quan đến thớ
tổng quát của nó, nên thừa số cơ sở nhất thiết phải là $K$.} trong đó $B$ là
một lược đồ abel trên $V$ (về căn bản là duy nhất, như ta đã nhắc).

Có thể giả sử $B'$ là thương của $A_K$ theo một lược đồ con nhóm $C'$.
(LƯU Ý: nói chung, $C'$ sẽ không «rút gọn», tức là các vành địa phương của nó
sẽ có những phần tử lũy linh.) Xét lược đồ con đóng $C$ của $A$ là «bao đóng»
của $C'$, tức là lược đồ con đóng nhỏ nhất của $A$ sao cho $C_K$ chứa $C'$.
Ta có $C_K=C'$, và từ việc $V$ là một vành định giá rời rạc, dễ dàng suy ra
rằng $C$ là một lược đồ con nhóm của $A$ trên $V$. Vì $A$ riêng trên
$\operatorname{Spec}(V)=S$, nên $C$ cũng vậy; mặt khác, $A$ thậm chí còn xạ
ảnh trên $S$. Do đó, ta có thể áp dụng (6.1) để dựng $A/C=B$,
\src{115}
đó là $B$ cần tìm.

Sau cùng, các lập luận về cơ bản đã biết cho phép suy ra từ (6.2) kết quả sau:

\result{ĐỊNH LÝ}{7.2} Cho $S$ là phổ của một vành Artin, $F$ và $G$ là các
lược đồ nhóm thuộc kiểu hữu hạn trên $S$, và $u:F\to G$ là một đồng cấu lược
đồ nhóm trên $S$. Giả sử:
\begin{enumerate}
\item[(i)] $F$ phẳng trên $S$;
\item[(ii)] hạt nhân của $u$ là hữu hạn.
\end{enumerate}
Trong những điều kiện này, lược đồ thương $G/F$ tồn tại; cấu xạ chính tắc
$G\to G/F$ là toàn ánh và mở, và các thớ của nó là các lớp tương đương theo
nghĩa tập hợp được xác định bởi các tác động bên phải của $F$ lên $G$. Sau
cùng, nếu $u$ là một đơn cấu, thì $G\to G/F$ là phẳng và cấu xạ
\[
  G\times F\longrightarrow G\times_{G/F}G
\]
là một đẳng cấu; nói cách khác, $G$ là một không gian thuần nhất chính trên
$G/F$, có nhóm cấu trúc (tác động bên phải) là $F$, hay đúng hơn là
$F\times_S(G/F)$ được xem như một lược đồ nhóm trên $G/F$ (xem [1], B,
\S~6).

\result{HỆ QUẢ}{7.3} Trong những điều kiện này, $G$ phẳng trên $S$ khi và chỉ
khi $G/F$ cũng vậy. Nếu điều kiện này được thỏa mãn, thì việc lấy thương
theo $F$ giao hoán với mọi phép mở rộng cơ sở của $S$; và nếu $F$ là một nhóm
con chuẩn tắc của $G$, thì $G/F$ có thể được trang bị một cấu trúc
\emph{nhóm thương} của $G$ theo $F$.

Tình hình đặc biệt đơn giản khi $S$ là phổ của một trường, vì khi đó mọi
$S$-tiền lược đồ đều đương nhiên phẳng trên $S$. Ta thu được:

\result{HỆ QUẢ}{7.4} Cho $F,G$ là các lược đồ nhóm thuộc kiểu hữu hạn trên một
trường $k$, và cho $u:F\to G$ là một đồng cấu $k$-nhóm. Khi đó $u$ phân tích
thành $F\to F'\to G$, trong đó $F\to F'$ là đồng cấu lấy thương theo nhóm con
đóng $\operatorname{Ker}u$ của $F$, còn $F'\to G$ là một đồng cấu nhóm và đồng
thời là một phép nhúng đóng. Lược đồ thương $G/F=G/F'$ tồn tại. Hình thức luận
thông thường (kiểu các định lý Noether) áp dụng được cho các nhóm đại số trên
$k$.

Kết quả này cho phép xử lý một cách thống nhất việc lấy thương trong các nhóm
đại số theo nghĩa cổ điển (tức là những nhóm bất khả quy trên $k$ và trơn trên
$k$), và việc lấy thương theo các nhóm con «vô cùng bé»
\src{116}
được CARTIER xét. Có lợi khi xem các «hyperalgebra» do CARTIER đưa ra, tiếp nối
các công trình của DIEUDONNÉ về các nhóm hình thức, như những nhóm trong phạm
trù các lược đồ hình thức trên $k$, và khi thích hợp (nếu chúng tương ứng với
các hyperalgebra có hạng hữu hạn trên $k$), như các nhóm đại số hữu hạn trên
$k$.

\fsection{8}{Một giả thuyết}

Nó đáp ứng nhu cầu lấy thương theo nhóm xạ ảnh tác động trên một số lược đồ con
của các «lược đồ Hilbert» (trong lý thuyết lược đồ, các lược đồ sau thay thế
các đa tạp Chow).

Cho $S$ là một tiền lược đồ và $n$ là một số nguyên. Với mỗi tiền lược đồ $S'$
trên $S$, ta cho tương ứng nhóm
$\operatorname{Gl}(n,\Gamma(S',\mathcal O_{S'}))$ gồm các ma trận khả nghịch có
$n$ hàng và $n$ cột, với các hệ số trong vành các tiết diện của
$\mathcal O_{S'}$. Ta thu được một hàm tử phản biến theo $S'$; dễ chứng minh
rằng nó biểu diễn được, nên tương ứng với một lược đồ nhóm trên $S$, hơn nữa
affine trên $S$, ký hiệu $\operatorname{Gl}(n)_S$. Phép dựng này tương thích
với đổi cơ sở, vì vậy trên thực tế tất cả xuất phát từ một lược đồ nhóm trên
$\mathbf Z$, ký hiệu $\operatorname{Gl}(n)$. Nhóm $\operatorname{Gl}(1)$, gọi
là nhóm nhân và thường ký hiệu $\mathbf G_m$, tương ứng với hàm tử
$S\longmapsto\Gamma(S,\mathcal O_S)^*$, nhóm các «đơn vị» trên $S$. Có một
đồng cấu hiển nhiên $\operatorname{Gl}(1)\to\operatorname{Gl}(n)$, và dễ dàng
dựng được nhóm thương, ký hiệu $\operatorname{GP}(n-1)$ và gọi là
\emph{nhóm xạ ảnh bậc $n-1$ trên $\mathbf Z$}. Nó biểu diễn hàm tử gán cho
$S$ nhóm các tiết diện của bó
$\widetilde{\operatorname{Gl}(n)}_S/
\widetilde{\operatorname{Gl}(1)}_S$, trong đó
$\widetilde{\operatorname{Gl}(n)}_S$ ký hiệu bó các mầm tiết diện của
$\operatorname{Gl}(n)_S$ trên $S$. (Cần lưu ý rằng các tiết diện của
$\operatorname{GP}(n-1)_S$ trên $S$ nói chung không xuất phát từ các tiết diện
của $\operatorname{Gl}(n)_S$ trên $S$!) Lưu ý rằng có thể chứng minh
$\operatorname{GP}(n-1)$ cũng biểu diễn hàm tử
$S\longmapsto\operatorname{Aut}_S(\mathbf P_S^{n-1})$, trong đó
$\mathbf P_S^{n-1}$ là không gian xạ ảnh chuẩn có chiều tương đối $n-1$ trên
$S$, ít nhất khi $S$ là Noether. Chính theo cách này mà nó xuất hiện trong lý
thuyết moduli.

Cho $S$ là một lược đồ Noether, nếu muốn có thể giả sử affine, và cho $X$ là
một $S$-tiền lược đồ tựa xạ ảnh, được trang bị một bó khả nghịch $\mathcal L$
rất rộng tương đối với $S$. Giả sử nhóm $G=\operatorname{GP}(n)_S$ tác động
trên $X$ và đồng thời trên $\mathcal L$ (tương thích với tác động của nó trên
$X$), và tác động \emph{tự do} trên $X$.\footnote{Văn bản in ghi «trên $S$».
Tác động đang xét và điều kiện lấy thương liên quan đến $X$.}
\src{117}

\result{GIẢ THUYẾT}{8.1}\footnote{Bản đính chính gắn với Chuyên đề 212 cho
biết giả thuyết này «hoàn toàn sai». Bản này cũng cho biết kết quả tích cực
trong Nhận xét 8.1 về sau đã được chứng minh và dẫn chiếu đến công trình của Mumford
cùng những cải tiến tiếp theo.} Trong các điều kiện trên:

\smallskip
\noindent $1^\circ$ \emph{Quan hệ tương đương do $G$ xác định là hữu hiệu,
thương $Y=X/G$ thuộc kiểu hữu hạn trên $S$, và cấu xạ chính tắc $f:X\to Y$ là
phẳng và toàn ánh} (do đó $X$ trở thành một không gian thuần nhất chính trên
$Y$, với nhóm $G\times_S Y=\operatorname{GP}(n)_Y$).

\smallskip
\noindent $2^\circ$ Cho $\mathcal L'$ là bó khả nghịch trên $Y$ được suy ra
từ $\mathcal L$ bằng «hạ xuống phẳng trung thành» theo $f$ (xem [1], B,
\S~th. 1). Khi đó $\mathcal L'$ là \emph{tiền rộng trên $Y$ tương đối với
$S$}, nghĩa là tồn tại một số nguyên $m$ và một cấu xạ tựa hữu hạn từ $Y$ vào
một không gian xạ ảnh chuẩn thích hợp $\mathbf P_S^N$, sao cho
$(\mathcal L')^{\otimes m}$ đẳng cấu với ảnh ngược của
$\mathcal O_{\mathbf P_S^N}(1)$.

Lưu ý rằng ngay cả khi $X$ tách trên $S$, vẫn có thể xảy ra việc $Y$ không
tách trên $S$ (một tình huống gặp trong những bài toán moduli hoàn toàn không
bệnh lý). Nếu $1^\circ$ đúng, thì $Y$ tách khi và chỉ khi quan hệ tương đương
do $G$ xác định có đồ thị đóng, nghĩa là khi $G\times X\to X\times X$ có ảnh
đóng (khi ấy nó là một phép nhúng đóng). Nếu $Y$ tách, thì $\mathcal L'$ tiền
rộng trên $Y$ tương đối với $S$ khi và chỉ khi nó rộng, nghĩa là khi một lũy
thừa tensor thích hợp xác định một phép nhúng xạ ảnh. Trong các bài toán moduli
được nhắc đến ở phần mở đầu, có thể chứng minh rằng quan hệ tương đương thu
được quả thật có đồ thị đóng.

\medskip
\noindent\textbf{NHẬN XÉT 8.1. ---} Ta đã giả sử
$G=\operatorname{GP}(n)_S$ để cố định ý tưởng và vì hiện nay đây là trường hợp
quan trọng nhất trong thực hành. Giả thiết hợp lý hơn đối với $G$ dường như là
$G$ thuộc một trong các «dạng» trên $S$ của một trong các nhóm Tohokû (phép
dựng các nhóm này trên các số nguyên cũng do CHEVALLEY thực hiện). Kết quả tích
cực duy nhất mà tôi biết theo hướng của giả thuyết trên là như sau: cho $X$ là
\emph{một lược đồ affine trên một trường $k$ có đặc trưng $0$, trên đó nhóm
$\operatorname{Gl}(n)_k$ hoặc $\operatorname{GP}(n-1)_k$ tác động tự do. Khi
đó quan hệ tương đương do $G$ xác định là hữu hiệu, thương $X/G$ cũng affine,
và cấu xạ $X\to X/G$ là phẳng và toàn ánh}. Chứng minh sử dụng sự kiện sau
(hiện chỉ được chứng minh trong đặc trưng $0$): khi cho $G$ tác động trên vành
tọa độ affine $A$ của $G$, xem như một không gian vectơ trên $k$, biểu diễn tầm
thường của $G$ chỉ xuất hiện một lần (trong một dãy hợp thành của một không
gian vectơ con hữu hạn chiều trên $k$, ổn định dưới $G$). Có vẻ như việc sử
dụng có hệ thống
\src{118}
lý thuyết biểu diễn tuyến tính của $G$ rốt cuộc sẽ cho một chứng minh của
giả thuyết, ít nhất khi cơ sở là một trường. Khi cơ sở không còn là một trường,
diễn giả không biết gì cả.

\section*{TÀI LIỆU THAM KHẢO}
\addcontentsline{toc}{section}{Tài liệu tham khảo}

\begin{description}
\item[{[1]}] GROTHENDIECK (Alexander). --- Techniques de descente et théorèmes
d'existence en géométrie algébrique, I: Généralités, Séminaire Bourbaki,
1959/60, n\textsuperscript{o} 190, 29 p.

\item[{[2]}] GROTHENDIECK (Alexander). --- Techniques de descente et théorèmes
d'existence en géométrie algébrique, II: Le théorème d'existence en théorie
formelle des modules, Séminaire Bourbaki, 1959/60,
n\textsuperscript{o} 195, 22 p.

\item[{[3]}] GROTHENDIECK (A.) và DIEUDONNÉ (J.). --- Éléments de géométrie
algébrique, I: Le langage des schémas. --- Paris, Presses universitaires de
France, 1960 (Institut des hautes Études scientifiques, Publications
mathématiques, 4).
\end{description}

\par
